\subsection{赋范空间的有限维向量子空间}

下面的命题将在 TA（待出版）中作为博苏克–乌拉姆定理的推论而证明。

定理 3. — 设 n 是正整数，V 是维数为 \(n + 1\) 的实赋范向量空间，W 是 V 的维数为 n 的向量子空间。设 S 是 V 的单位球面。不存在连续映射 \(f: S \to W\) ，使得 \(\|f(x) - x\| < 1\) 对一切 \(x \in S\) 成立。

定理 4（克赖因、克拉索塞利斯基、米尔曼）. — 设 E 是赋范空间，F 和 G 是 E 的向量子空间。假设 G 的维数有限且严格小于 F 的维数。则存在 F 中范数为 1 的元素，它到 G 的距离等于 1。

只需处理域 K 等于 R 的情形。设 n 是 G 的维数。把 F 换成 F 的一个包含 G 且维数为 \(n+1\) 的向量子空间，就化归于 \(\dim(F)=n+1\) 的情形。用反证法，假设定理的结论不成立。设 S 是 F 的单位球面。对每个 \(x\in S\) ，x 到 G 的距离严格小于 \(\|x\|=1\) ，故可选取 G 的元素 \(y(x)\) ，使得 \(\|x-y(x)\|<1\) 。用 \(V_{x}\) 表示 S 中满足 \(\|z-y(x)\|<1\) 的元素 z 所成的集；它是 x 在 S 中的一个开邻域。存在 S 上的局部有限连续单位分解 \((\varphi_{x})_{x\in\mathrm{S}}\) ，它从属于 S 的覆盖 \((\mathrm{V}_{x})_{x\in\mathrm{S}}\) （TG, IX, p. 46, 命题 3 及 p. 51, 命题 6）。用 \(f\colon S\to G\) 表示由下式给出的连续映射

\[
f (z) = \sum_ {x \in \mathrm{S}} \varphi_ {x} (z) y (x).
\]

对一切 \(z \in S\)，我们有 \(\varphi_x(z) \geqslant 0\) （对一切 \(x \in S\)），且存在 \(x \in S\) 使 \(\varphi_x(z) > 0\)，因为 \(\sum_{x \in S} \varphi_x(z) = 1\)，由此

\[
\begin{array}{c} \| z - f (z) \| = \left\| \sum_ {x \in \mathrm{S}} \varphi_ {x} (z) (z - y (x)) \right\| \leqslant \sum_ {x \in \mathrm{S}} \varphi_ {x} (z) \| z - y (x) \| \\ <   \sum_ {x \in \mathrm{S}} \varphi_ {x} (z) = 1. \end{array}
\]

f 的这一性质与定理 3 矛盾。

命题 9. — 设 E 和 F 是赋范空间，\(n \in N\)，并设 u 是从 E 到 F 中的连续线性映射，其核的维数为 n。u 的余范数等于从 u 到映射 \(v \in \mathcal{L}(\mathrm{E};\mathrm{F})\) （其核的维数至少为 \(n + 1\)）所成的集的距离 d。

当 \(u = 0\) 时，我们有 \(((u)) = +\infty\) 且 \(\dim(\mathrm{E}) = n\) ，从而 \(d = +\infty\) 。以下设 \(u\) 非零，于是 \(((u)) < +\infty\) 。设 \(v\) 是 \(\mathcal{L}(\mathrm{E};\mathrm{F})\) 中满足 \(\| u - v \| < ((u))\) 的元素，我们来证明它的核的维数 \(\leqslant n\) 。设 \(x\) 是 \(\operatorname{Ker}(v)\) 的一个范数为 1 的元素，\(y\) 是它在 \(\operatorname{E}/\operatorname{Ker}(u)\) 中的像。我们有（III, p. 61 的公式 (5)）

\[
((u)) \| y \| \leqslant \| u (x) \| = \| (u - v) (x) \| \leqslant \| u - v \| <   ((u))
\]

从而 \(\|y\| < 1\) 。而 \(\|y\|\) 正是 x 到 \(\operatorname{Ker}(u)\) 的距离。于是定理 4 蕴含 \(\dim \operatorname{Ker}(v) \leqslant \dim \operatorname{Ker}(u) = n\) 。这就证明了不等式 \(((u)) \leqslant d\) 。反向不等式 \(d \leqslant ((u))\) 由下面更精确的引理得出。

引理 3. — 设 c 是满足 \(c > ((u))\) 的实数。存在连续线性映射 \(h: E \to F\) ，其秩为 1、范数 \textless{} c，使得 \(u + h\) 的核包含 u 的核且维数为 \(n + 1\) 。

用 \(p\) 表示从 E 到 \(\mathrm{E} / \mathrm{Ker}(u)\) 上的典范映射。存在 \(a \in \mathrm{E}\) ，使得 \(\| p(a) \| = 1\) 且 \(\| u(a) \| < c\) （III, p. 61 的公式 (6)）。由哈恩–巴拿赫定理（EVT, II, p. 24, 推论 2），存在连续线性泛函 \(f\) （在赋范空间 \(\mathrm{E} / \mathrm{Ker}(u)\) 上），使得 \(f(p(a)) = 1\) 且 \(\| f \| = 1\) 。线性映射 \(h: x \mapsto -(f \circ p)(x)u(a)\) （从 E 到 F 中）连续且秩为 1，并且有 \(\| h \| \leqslant \| f \| \| p \| \| u(a) \| < c\) 。映射 \(u + h\) 的核包含 \(u\) 的核和 \(a\) ；由于 \(a \notin \mathrm{Ker}(u)\) ，它的维数至少为 \(n + 1\) 。另一方面，线性映射 \(u\) 通过过渡到子空间诱导一个从 \(\mathrm{Ker}(u + h)\) 到 \(\mathrm{Im}(h)\) 中的线性映射，其核为 \(\mathrm{Ker}(u) \cap \mathrm{Ker}(u + h)\) ，于是 \(\dim(\mathrm{Ker}(u + h)) \leqslant \dim(\mathrm{Ker}(u)) + 1\) 。结论得证。

推论 1. --- 设 E 和 F 是赋范空间，u、v 是从 E 到 F 中的非零连续线性映射，其核有相同的有限维数。则我们有

\[
\left| \left((u)\right) - ((v)) \right| \leqslant \| u - v \|.
\]

用 \(n\) 表示 \(u\) 与 \(v\) 的核的共同维数。用 A 表示从 E 到 F 的、核的维数 \(\geqslant n + 1\) 的连续线性映射所成的集。由于 \(u \neq 0\) ，我们有 \(\dim(\mathrm{E}) > n\) ，且集 A 包含 0。由命题 9，\(((u))\) 与 \(((v))\) 分别是 \(u\) 与 \(v\) 到 \(\mathcal{L}(\mathrm{E};\mathrm{F})\) 中集 A 的距离。于是只需应用公式 \(|d(u,\mathrm{A}) - d(v,\mathrm{A})| \leqslant \|u - v\|\) （参见 TG, IX, p. 13）。

推论 2. --- 设 E 和 F 是巴拿赫空间，u 和 v 是从 E 到 F 中的非零连续线性映射，其像在 F 中有相同的有限余维数。则有

\[
| ((u)) - ((v)) | \leqslant \| u - v \|.
\]

态射 u 是严格的（III, p. 52, 引理 6）。连续线性映射 \({}^{t}u\) 的核是正交 \((\mathrm{Im}(u))^{\circ}\) （\(\mathrm{Im}(u)\) 的正交）（EVT, IV, p. 27, 命题 2），由此 \(\dim(\mathrm{Ker}({}^{t}u)) = \mathrm{codim}_{\mathrm{F}}(\mathrm{Im}(u))\) ，类似地 \(\dim(\mathrm{Ker}({}^{t}v)) = \mathrm{codim}_{\mathrm{F}}(\mathrm{Im}(v))\) 。因此 \({}^{t}u\) 与 \({}^{t}v\) 的核有相同的有限维数。由于

\[
\| ^ {t} u - ^ {t} v \| = \| u - v \|, \quad ((^ {t} u)) = ((u)), \quad ((^ {t} v)) = ((v))
\]

（EVT, IV, p. 7, 命题 8 及 III, p. 63, 命题 8），把推论 1 应用于 \(^{t}u\) 和 \(^{t}v\) 即得该断言。

